% id: 33-27
% book: 개념원리 중학 3-1 · 03 무리수와 실수
% section: 중단원 마무리하기 STEP 3 실력 UP
% figure: none
% answer: $54$
% answer_source: 답지
% variant_level: 0 (원본 전사)
% 이 파일은 build.mjs 가 items.json 에서 만든다. 고칠 때는 items.json 을 고친다.
\smash{\raisebox{3.5mm}{\dmkichul{개념원리 중학 3-1 33쪽 27번}}}
자연수 $x$에 대하여 $\sqrt{x}$ 이하의 자연수의 개수를 $N(x)$라 하자. 예를 들어 $2<\sqrt{5}<3$이므로 $N(5)=2$이다. 이때 $N(1)+N(2)+N(3)+\cdots+N(20)$의 값을 구하시오.
